\section{Introduction}\label{sec:intro}
\subsection{Background}
An anonymous credential lets a holder prove a statement about themselves while revealing nothing else. Education is one place they are starting to be used. A school issues a credential, and a student uses it to prove a single fact, such as being enrolled or eligible for a programme, without revealing the rest of their record. Forged credentials cause real harm. Falsified admission letters were behind recent student removal proceedings~\cite{bbc2023canada}. The data a credential carries is also a target. The $2021$ scraping of some $700$ million LinkedIn profiles is one example~\cite{nixon2021linkedin}. This threat is also growing over time. Standards bodies now tell organisations to migrate to post-quantum cryptography~\cite{nistir8547}, and industry is following. Google has committed to finish its own migration by $2029$~\cite{adkins2026google}. A credential issued today must therefore be three things at once: unforgeable, anonymous, and secure against a future quantum computer. A recent construction, BDEC~\cite{10.1007/978-981-96-0957-4_3}, shows the three can be had together, by combining a post-quantum signature with a post-quantum-secure proving system. Showing a credential means proving you hold a valid signature without revealing it, and the proving system is what makes that possible.

Building the credential once, however, is not the end. One meant to last a decade must keep working as two things around it change. First, what the holder must prove can change. A verifier may later ask for something the credential did not prove before. Second, the signature the credential is built on may be replaced as post-quantum standards mature. In a zkSNARK, what will be proved must be decided before the proof system is built, so either change forces the whole verification to be rebuilt and re-checked. A zkVM (zero-knowledge virtual machine) instead takes the computation as an input, so either change is only an edit to the program it runs.

\subsection{From Constructing to Deploying}\label{ssec:deploy-question}
This flexibility has a cost. A zkVM can run any program, but whenever a program computes on numbers larger than the prover handles natively, it pays for it, and the signatures we care about do exactly that. BDEC's efficiency rests on the signature scheme Loquat~\cite{cryptoeprint:2024/868}, for which this thesis substitutes its successor PLUM~\cite{10.1007/978-981-95-2961-2_6}. PLUM reports smaller signatures and faster verification at the same security level (Section~\ref{sec:related}). PLUM computes on very large numbers ($192$-bit, and a concrete $199$-bit prime as instantiated; Section~\ref{sec:prelim}), while a zkVM works natively only with numbers of $31$ bits. Each operation of the signature must therefore be emulated with many small-number operations, at high cost. Section~\ref{sec:prelim} names this cost the \emph{field-mismatch tax}. It is a real barrier. On the target hardware, verifying a single signature this way ran for hours and never produced a proof (Section~\ref{sec:evaluation}).

This is the tension the thesis addresses. It asks what is required to run PLUM inside a zkVM on the consumer hardware a student or an institution would actually own, and at what cost. The accounting includes the changes that motivate the zkVM in the first place, both to what the credential proves and to the signature itself.

%%% Perhaps the state-of-art could go here as a sentence that links back towards Related Works? %%%

\subsection{The Precompile Suite, and the Question It Answers}\label{ssec:question}
The zkVM measured here is SP1~\cite{succinct2024sp1}, chosen because it lets the builder add \emph{application-defined precompiles}. A precompile is a custom accelerator for one operation, which the program calls directly, so the operation runs at full speed instead of being emulated. The more work a proof performs, the more it costs. A precompile removes the field-mismatch tax for an operation that would otherwise dominate that cost, in exchange for the smaller cost of its own accelerator. This thesis builds a suite of precompiles for PLUM verification and uses it as its instrument of study. Building the suite turns a verification that never finishes into one that completes, with a cost that can be measured and explained. It also forces a sharper question. Is a custom precompile worth building? The answer depends on what you compare it against. Without a precompile the hash never finishes, so the precompile for PLUM's hash, Griffin, is what lets PLUM run at all. Yet a plain software SHA-3, which needs no precompile, proves \emph{faster} than the accelerated Griffin (Section~\ref{sec:evaluation}). A precompile can therefore be necessary to make an operation run and still not be worth building, when a simpler alternative proves faster. The SHA-3 version is a diagnostic control, not a deployable form of PLUM. This smaller question sits inside the central question of this thesis:
%is nested inside the central question of this thesis:

\begin{quote}
\emph{Can a post-quantum anonymous credential run inside a zkVM on consumer hardware without breaking the security properties it must provide, and when plain execution cannot finish, what makes it run?}
\end{quote}

The question separates into two layers. The feasibility layer is priced by a cost model (Section~\ref{sec:theoretical}) that weighs the cost a precompile adds against the cost it removes. The security layer asks which of the three requirements survive the move, and each becomes a property the proof must have. Unforgeability needs the proof to be \emph{knowledge-sound}, meaning only someone who genuinely holds a valid signature can produce a convincing proof. The zkVM provides this, so unforgeability transfers, under the assumptions of Section~\ref{sec:security}. Anonymity needs the proof to be \emph{zero-knowledge}, revealing nothing beyond the fact it proves. Quantum-safety needs the proof to stay sound against an attacker with a quantum computer. \emph{Deployable post-quantum anonymity} (Definition~\ref{def:deployable}) asks for all three at once, in a single proof within the memory budget, and it is the one requirement SP1 cannot currently meet. Its default proof is quantum-safe but not zero-knowledge. The only proof its toolchain adds to supply zero-knowledge is pairing-based (Section~\ref{sec:prelim}), and it rests on elliptic-curve cryptography that a quantum computer breaks. So on the route SP1 ships today, a single proof is quantum-safe or anonymous, never both. This is a limit of the primitives available today, not a fact about zero-knowledge virtual machines. A fixed circuit reaches both at once, but only by giving up the \emph{flexibility} the zkVM was chosen for. The deployment choice thus trades flexibility against post-quantum anonymity. The dual obstruction is formalised in Section~\ref{sec:security}.

The credential construction (BDEC~\cite{10.1007/978-981-96-0957-4_3}) instantiated with PLUM is the concrete workload through which the cost criterion built on that model (Section~\ref{sec:theoretical}) is calibrated and measured on SP1. The criterion decides one question, whether a precompile is worth building. The choice of substrate is a second decision, weighed by the substrate-selection rule of Section~\ref{sec:discussion}. Both belong to the operator of a real deployment, such as a university consortium choosing a proving substrate for digital diplomas today. Such an operator must commit to a substrate before knowing which predicates its verifiers will demand or which parameter sets the standards process will finalise. The operator needs to know both what the precompile delivers and where it stops.

% [two-provers tikz figure unchanged; omitted here for brevity — keep yours]

\subsection{Contributions of This Thesis}
This thesis analyses which of the credential's security properties survive the move into a zkVM, locating the one that does not, deployable post-quantum anonymity; distils a reusable cost criterion for when a precompile pays; builds the precompile suite that turns an emulation that does not complete within the budget into a measurable proof; and reports the first empirical characterisation of PLUM verification inside a zkVM. The contributions are as follows.

\begin{enumerate}
  \item \textbf{The dual obstruction to post-quantum anonymity, and the everlasting-anonymity boundary.} The credential's unforgeability transfers to the substrate, because it reduces to the zkVM's knowledge soundness, which the zkVM provides. Deployable post-quantum anonymity does not transfer, and the obstruction has two sides. SP1's default proof (the object its toolchain calls a \emph{receipt}) is feasible and post-quantum, but it is not zero-knowledge (the first obstruction). The zero-knowledge wrap SP1's toolchain adds does fit the consumer memory budget, but it is pairing-based, resting on elliptic-curve cryptography that quantum computers break, so it is not post-quantum (the second obstruction). Feasibility is therefore not the obstruction; a limit of the available primitives is. A deployer who accepts classical security can close the first obstruction with the classical wrap, which leaves the second as the one that survives. Pushing the question to its limit, we prove that \emph{everlasting} anonymity, safety even against an attacker with unlimited computing time after the fact, cannot be established in the standard model (that is, without idealised oracles) for the class of hash-committed showings this construction uses, and we give the conditional route out (Section~\ref{sec:security}). The obstruction is formalised in Section~\ref{sec:security} and measured in Section~\ref{sec:evaluation}.

  \item \textbf{A cost criterion for when a precompile pays, and the hash-cost inversion it predicts.} Whether a precompile \emph{pays} is really two questions the one word runs together. Against doing the operation without a precompile, it pays when the cost it removes outweighs the cost of its own accelerator, a feasibility test. Against the best alternative for the same operation, it pays only if it adds less total proving cost than that alternative. Our cost model (Section~\ref{sec:theoretical}) answers both from a precompile's size, before it is built. The Griffin hash precompile over PLUM's field ($\mathbb{F}_p^{192}$)~\cite{cryptoeprint:2022/403} passes the first test and fails the second. It recovers a proof that emulation cannot finish, yet proves \emph{slower} than a software SHA-3 control, because its own accelerator adds proving cost a plain hash never pays. That comparison is run at $\lambda{=}80$, below PLUM's headline security level, because the full level does not complete on the target hardware. The control is matched on security level, not on field or hash, and is a diagnostic, not a deployable variant (Section~\ref{sec:evaluation}). The model separates the three precompiles the thesis ships the same way. Griffin pays for feasibility but loses to the SHA-3 control. The field-multiplication precompile is predicted not to pay, because SP1 already ships a multiplication precompile (its \texttt{UINT256\_MUL} accelerator) that serves it (Section~\ref{sec:system}). The power-residue check, which unlike the hash cannot be moved to a smaller field (Section~\ref{sec:system}), still pays on total proving cost (Section~\ref{sec:evaluation}). The criterion is reusable within this class. For another algebraic primitive on a small-field zkVM, a builder can decide from the specification alone whether a precompile lowers total proving cost, and the model gives a testable matched-field forecast, checked in Section~\ref{sec:evaluation}.

  \item \textbf{A precompile suite for PLUM verification over a $199$-bit prime.} We build three precompiles over the same field, one for the Griffin hash, one for field multiplication, and one for the power-residue check, each with a written argument that its accelerator computes the same function as the reference specification (this argument is partially discharged at the lowest level; Section~\ref{sec:security}). All three are built and test-proved, but only the Griffin precompile is on the measured path (Section~\ref{sec:system}). The Griffin precompile is what recovers feasibility.

  \item \textbf{The first empirical characterisation of PLUM-in-a-zkVM.} We write the credential's issuing and showing steps as SP1 programs with PLUM in place of Loquat, and run PLUM verification on its own, measuring both how many steps it takes and how long it takes to prove (Section~\ref{sec:evaluation}, where every figure appears with its full configuration). Without a precompile, verification was stopped unfinished after five hours. With the Griffin precompile, five proofs averaged $14.24$ minutes at $\lambda=80$. To our knowledge, this is the first such measurement of PLUM verification inside a zkVM.
\end{enumerate}

\noindent\emph{Security scope.} We give one conditional theorem that the security is preserved when PLUM runs inside the zkVM (Section~\ref{sec:security}). In the classical random-oracle model, it bounds how much an attacker can gain against the credential system's unforgeability, anonymity, and unlinkability. Anonymity and unlinkability are the two parts of the anonymity requirement from Section~\ref{sec:intro}. Whether anonymity survives a quantum attacker is the subject of the dual obstruction. The positive results rest on five explicitly named open assumptions. They also rest on one named conjecture (an extractor conjecture), needed because a long run's proof is split into pieces and recursively combined, and the argument that this combined proof is knowledge-sound is not yet complete (Section~\ref{sec:prelim}, Section~\ref{sec:security}).

\subsection{Thesis Organisation}
The thesis proceeds as follows: related work (Section~\ref{sec:related}) and primitives (Section~\ref{sec:prelim}); the cost model and update churn at both levels (Section~\ref{sec:theoretical}); the construction and precompile suite (Section~\ref{sec:system}); the security-preservation theorem and its open assumptions (Section~\ref{sec:security}); the measurements (Section~\ref{sec:evaluation}); the deployment decision framework (Section~\ref{sec:discussion}); and the conclusion (Section~\ref{sec:conclusion}).